\documentclass[10pt, letterpaper]{article}

% Packages:
\usepackage[
    ignoreheadfoot, % set margins without considering header and footer
    top=1 cm, % seperation between body and page edge from the top
    bottom=1 cm, % seperation between body and page edge from the bottom
    left=1 cm, % seperation between body and page edge from the left
    right=1 cm, % seperation between body and page edge from the right
    footskip=1 cm, % seperation between body and footer
    % showframe % for debugging
]{geometry} % for adjusting page geometry
\usepackage{titlesec} % for customizing section titles
\usepackage{tabularx} % for making tables with fixed width columns
\usepackage{array} % tabularx requires this
\usepackage[dvipsnames]{xcolor} % for coloring text
\definecolor{primaryColor}{RGB}{0, 0, 0} % define primary color
\usepackage{enumitem} % for customizing lists
\usepackage{fontawesome5} % for using icons
\usepackage{amsmath} % for math
\usepackage[
    pdftitle={John Doe's CV},
    pdfauthor={John Doe},
    pdfcreator={LaTeX with RenderCV},
    colorlinks=true,
    urlcolor=primaryColor
]{hyperref} % for links, metadata and bookmarks
\usepackage[pscoord]{eso-pic} % for floating text on the page
\usepackage{calc} % for calculating lengths
\usepackage{bookmark} % for bookmarks
\usepackage{lastpage} % for getting the total number of pages
\usepackage{changepage} % for one column entries (adjustwidth environment)
\usepackage{paracol} % for two and three column entries
\usepackage{ifthen} % for conditional statements
\usepackage{needspace} % for avoiding page brake right after the section title
\usepackage{iftex} % check if engine is pdflatex, xetex or luatex

% Ensure that generate pdf is machine readable/ATS parsable:
\ifPDFTeX
    \input{glyphtounicode}
    \pdfgentounicode=1
    \usepackage[T1]{fontenc}
    \usepackage[utf8]{inputenc}
    \usepackage{lmodern}
\fi

\usepackage{charter}

% Some settings:
\raggedright
\AtBeginEnvironment{adjustwidth}{\partopsep0pt} % remove space before adjustwidth environment
\pagestyle{empty} % no header or footer
\setcounter{secnumdepth}{0} % no section numbering
\setlength{\parindent}{0pt} % no indentation
\setlength{\topskip}{0pt} % no top skip
\setlength{\columnsep}{0.15cm} % set column seperation
\pagenumbering{gobble} % no page numbering

\titleformat{\section}{\needspace{4\baselineskip}\bfseries\large}{}{0pt}{}[\vspace{1pt}\titlerule]

\titlespacing{\section}{
    % left space:
    -1pt
}{
    % top space:
    0.3 cm
}{
    % bottom space:
    0.2 cm
} % section title spacing

\renewcommand\labelitemi{$\vcenter{\hbox{\small$\bullet$}}$} % custom bullet points
\newenvironment{highlights}{
    \begin{itemize}[
        topsep=0.10 cm,
        parsep=0.10 cm,
        partopsep=0pt,
        itemsep=0pt,
        leftmargin=0 cm + 10pt
    ]
}{
    \end{itemize}
} % new environment for highlights


\newenvironment{highlightsforbulletentries}{
    \begin{itemize}[
        topsep=0.10 cm,
        parsep=0.10 cm,
        partopsep=0pt,
        itemsep=0pt,
        leftmargin=10pt
    ]
}{
    \end{itemize}
} % new environment for highlights for bullet entries

\newenvironment{onecolentry}{
    \begin{adjustwidth}{
        0 cm + 0.00001 cm
    }{
        0 cm + 0.00001 cm
    }
}{
    \end{adjustwidth}
} % new environment for one column entries

\newenvironment{twocolentry}[2][]{
    \onecolentry
    \def\secondColumn{#2}
    \setcolumnwidth{\fill, 4.5 cm}
    \begin{paracol}{2}
}{
    \switchcolumn \raggedleft \secondColumn
    \end{paracol}
    \endonecolentry
} % new environment for two column entries

\newenvironment{threecolentry}[3][]{
    \onecolentry
    \def\thirdColumn{#3}
    \setcolumnwidth{, \fill, 4.5 cm}
    \begin{paracol}{3}
    {\raggedright #2} \switchcolumn
}{
    \switchcolumn \raggedleft \thirdColumn
    \end{paracol}
    \endonecolentry
} % new environment for three column entries

\newenvironment{header}{
    \setlength{\topsep}{0pt}\par\kern\topsep\centering\linespread{1.5}
}{
    \par\kern\topsep
} % new environment for the header

\newcommand{\placelastupdatedtext}{% \placetextbox{<horizontal pos>}{<vertical pos>}{<stuff>}
  \AddToShipoutPictureFG*{% Add <stuff> to current page foreground
    \put(
        \LenToUnit{\paperwidth-2 cm-0 cm+0.05cm},
        \LenToUnit{\paperheight-1.0 cm}
    ){\vtop{{\null}\makebox[0pt][c]{
        \small\color{gray}\textit{Last updated in September 2024}\hspace{\widthof{Last updated in September 2024}}
    }}}%
  }%
}%

% save the original href command in a new command:
\let\hrefWithoutArrow\href

% new command for external links:


\begin{document}
    \newcommand{\AND}{\unskip
        \cleaders\copy\ANDbox\hskip\wd\ANDbox
        \ignorespaces
    }
    \newsavebox\ANDbox
    \sbox\ANDbox{$|$}

    \begin{header}
        \fontsize{25 pt}{25 pt}\selectfont Avishek Bose \\
        \fontsize{10 pt}{10 pt}\selectfont AI/ML Researcher

        \normalsize
        Location: Knoxville, TN, United States | Immigration status: Permanent resident (US)
        \\
        \mbox{\hrefWithoutArrow{mailto:avishek85.15th@gmail.com}{avishek85.15th@gmail.com}}%
        \kern 5.0 pt%
        \AND%
        \kern 5.0 pt%
        \mbox{\hrefWithoutArrow{+1(785)317-6675}{(785)317-6675}}%
        \kern 5.0 pt%
        \AND%
        \kern 5.0 pt%
        \mbox{\hrefWithoutArrow{https://linkedin.com/in/avibose}{linkedin.com/in/avibose}}%
        \kern 5.0 pt%
    \end{header}

    \vspace{- 0.3 cm}
    \section{Summary}
    AI/ML researcher with 8+ years of experience developing deep learning systems for scientific applications. Published author of 3 journal articles and 22 peer-reviewed conference papers in venues including IPDPS, IEEE BigData, and ICMLA. Expertise in GNNs and Transformer architectures, with direct experience applying these methods to drug toxicology prediction, scientific foundation model pretraining, and large-scale knowledge extraction from domain-specific corpora.

    \vspace{-1mm}
    \section{Work Experience}

    \begin{twocolentry}{Jun 2025 -- Present}
        \textbf{AI/ML Platform Engineer}
    \end{twocolentry}
    \begin{twocolentry}{Knoxville, TN, USA}
       Tegna Inc.
    \end{twocolentry}

    \begin{onecolentry}
        \begin{highlights}
            \item Designed and deployed a large-scale RAG platform on Azure serving 100K+ daily production requests, integrating semantic search, summarization, and LLM-based recommendation workflows.
            \item Established automated evaluation pipelines using RAGAs and DeepEval to continuously monitor retrieval precision and generation quality across 5+ production models.
        \end{highlights}
    \end{onecolentry}

\vspace{1mm}

    \begin{twocolentry}{Jan 2023 -- Jun 2025}
        \textbf{Postdoctoral Research Associate}
    \end{twocolentry}
    \begin{twocolentry}{Oak Ridge, TN, USA}
        Critical Infrastructure Resilience group, Geospatial Science and Human Security Division, Oak Ridge National Laboratory (ORNL)
    \end{twocolentry}
    \vspace{0.10 cm}
    \begin{onecolentry}
        \begin{highlights}
            \item Applied GNN and Transformer-based architectures to drug toxicology prediction and materials property inference, translating graph-structured molecular representations into classification and regression tasks.
            \item Co-developed MatGPT, a scientific foundation model pretrained on 2M+ materials science documents using Frontier (world's first Exascale supercomputer), demonstrating scalable pretraining on domain-specific scientific corpora.
            \item Built LLM fine-tuning pipelines using LoRA for domain adaptation across 50K+ DOE/ORNL reports, achieving $\sim$15\% accuracy improvement over zero-shot baselines.
            \item Designed cloud-native ML inference infrastructure on Azure supporting 24$\times$7 scientific workloads with $<$100ms response time and 99.9\%+ availability.
        \end{highlights}
    \end{onecolentry}

    \vspace{1mm}

    \begin{twocolentry}{Jun 2018 -- Dec 2022}
        \textbf{Research Assistant}
    \end{twocolentry}
    \begin{twocolentry}{Manhattan, Kansas, USA}
       KDD Lab, Department of Computer Science, Kansas State University
    \end{twocolentry}
    \vspace{0.10 cm}
    \begin{onecolentry}
        \begin{highlights}
            \item Developed GNN-based models achieving $\sim$89\% accuracy for predicting HPC job execution status and resource requirements from structured cluster logs.
            \item Built NLP and graph-based systems to classify and extract cyber-threat intelligence events from 10M+ social media posts.
            \item Designed abstractive summarization pipelines for disaster situational reports, processing 1K+ reports per batch using generative sequence-to-sequence models.
        \end{highlights}
    \end{onecolentry}

    \vspace{0.1 cm}
    \begin{twocolentry}{Jan 2015 -- Dec 2016}
        \textbf{Software Engineer}
    \end{twocolentry}
    \begin{twocolentry}{Dhaka, Bangladesh}
       ICT Division, Dutch Bangla Bank
    \end{twocolentry}

    \begin{onecolentry}
        \begin{highlights}
            \item Developed backend data processing and reporting systems handling 100K+ daily banking transactions using SQL and enterprise software platforms.
        \end{highlights}
    \end{onecolentry}

    \section{Publications}
        3 journal and 22 conference papers in IPDPS, BigData, ICMLA, ISI \\
        \hrefWithoutArrow{https://scholar.google.com/citations?user=y5ndUuMAAAAJ\&hl=en\&oi=ao}{Google Scholar: scholar.google.com/citations?user=y5ndUuMAAAAJ}

\section{Education}

        \begin{twocolentry}{
            August 2017 -- Dec 2022
        }
            \textbf{Kansas State University}, Ph.D. in Computer Science\end{twocolentry}

        \begin{twocolentry}{
            January 2014 -- Dec 2015
        }
            \textbf{University of Dhaka}, M.S. in Computer Science and Engineering\end{twocolentry}

        \begin{twocolentry}{
            January 2009 -- Aug 2013
        }
            \textbf{University of Dhaka}, B.Sc. in Computer Science and Engineering\end{twocolentry}

\vspace{-1mm}
\section{Technical Skills}

        \paragraph{AI/ML frameworks and tools:} PyTorch, Hugging Face, PyTorch Geometric, Scikit-learn, LlamaIndex
        \vspace{-4mm}
        \paragraph{AI/ML:} GNN, Transformer fine-tuning, Foundation model pretraining, LLM fine-tuning, RAG, VectorDB, distillation, quantization, prompt tuning
        \vspace{-4mm}
        \paragraph{Cloud \& MLOps:} Azure, MLflow, Docker, Terraform, CI/CD
        \vspace{-4mm}
        \paragraph{Programming Languages:} Python, C\#, C

\end{document}
